\documentclass[11pt]{article}
\usepackage[margin=0.85in]{geometry}
\usepackage[T1]{fontenc}
\usepackage{amsmath,amssymb,booktabs,graphicx,xurl,hyperref}
\graphicspath{{outputs/task1/}{outputs/task2/}{outputs/task3/}}

\title{CS 498 HW0 Report\\Perspective, Homography, and Epipolar Geometry}
\author{Your Name (NetID)}
\date{Fall 2026}

\begin{document}
\maketitle

\section{Homography and 2D insertion}
State your point normalization and DLT construction. Explain the transform
from logo pixels to the metric court rectangle and then to the image. Include
\path{task1a_basketball_homography.png} and
\path{task1b_basketball_insertion.png}, report the maximum corner-transfer
error, and explain your alpha-composition equation.

\subsection{Tennis transfer and camera height}
List the A--H metric XY coordinates you derived for the supplied UV points. Include
\path{task1c_tennis_correspondences.png},
\path{task1c_tennis_insertion.png} and
\path{task1c_tennis_alignment.png}. Report the eight-anchor reprojection RMSE and explain why
a planar homography with unknown intrinsics does not determine a unique camera
height.

\section{Camera projection and 3D insertion}
Describe the projection DLT and explain how knowledge of $K$ determines whether
the metric court plane is sufficient or non-coplanar landmarks are needed.
Include \path{task2b_bunny_insertion.png},
\path{task2a_reprojection.png}, and \path{task2c_camera_pose.png}. Report
the mean and maximum reprojection error, $K$, $R$, $t$, $C$, and the optical-center
height. Verify $C=-R^\top t$.

\section{Two-view epipolar geometry}
Describe the two point normalizations, eight-point design matrix, and rank-two
enforcement. Include \path{task3b_correspondences.png} and
\path{task3b_epipolar_lines.png}. Report the rank, mean symmetric epipolar
distance, and median distance; explain why real correspondence residuals are
not exactly zero.

\section{Limitations and improvements}
Discuss annotation precision, degenerate configurations, and at least two
possible improvements.

\section*{Optional foreground-preserving video bonus}
If attempted, describe your background estimate, foreground alpha, and final
composition. State that you implemented the complete bonus pipeline without a
provided helper. Submit \path{bonus_tennis_video.py}, one representative PNG
for each of match001, match005, match028, match091, and match125, and one final
MP4 for match005. This is the only section in which Gen-AI leverage is allowed.

\section*{AI-use declaration}
Check one statement.

\begin{itemize}
  \item[$\square$] I did not use generative AI for this homework.
  \item[$\square$] I did not use generative AI for any required part. I used it only in the optional bonus marked ``Gen-AI leverage allowed.'' Tool/model, purpose, and affected files: \underline{\hspace{1.4in}}
\end{itemize}

I confirm that generative AI did not write or complete my required code,
select required correspondences, derive my required answers, generate my
analysis, or write my required report.

\end{document}
